\subsection{On the Efficient Implementation of Cryptographic Algorithms}\label{sec:security_of_cryptographic_algorithms:performance}


%\remarkBlock{The many declinations of efficiency}{Efficiency may refer to several aspects of software development and execution, such as performance (measured in time or processor cycles), memory efficiency (measures in bytes stored or number of memory accesses or cache hit/miss ratio), power consumption (measured in Watts or Joules), cost-effectiveness (measure monetarily), and complexity efficiency (measured in expertise and skill required to develop software). In our context, we consider efficiency in terms of performance and memory efficiency.}

%At the same time, efficiency is critical to meet the requirements of real-world applications, ranging from high-speed financial transactions to resource-constrained devices. In this regard, the approach taken in the implementation of cryptographic algorithms greatly influences performance. Finally, it may not always be easy to balance security and efficiency in the implementation of cryptographic algorithms due to their interplay: strong security measures must be evaluated against their computational and operational overhead, while optimized implementations must be analyzed for potential weaknesses. This interplay underscores the importance of systematically ensuring and evaluating both the security and efficiency of cryptographic implementations to build resilient and performant systems.

Below, we discuss efficiency of cryptographic algorithms from software implementation (i.e., program) to hardware execution (i.e., process). In fact, although the efficiency of a program is usually measured asymptotically with respect to the given input (i.e., computational complexity), its actual efficiency ultimately depends on its implementation. First, we present the most relevant best practices for optimizing cryptographic algorithm implementation (\Cref{sec:security_of_cryptographic_algorithms:performance:best_practices}) and compare  different programming languages with respect to their efficiency and the optimizations available in the corresponding compilers and interpreters (\Cref{sec:security_of_cryptographic_algorithms:performance:languages_and_compilers}). Then, we list available hardware-based technologies for improving the efficiency of the execution of cryptographic algorithms (\Cref{sec:security_of_cryptographic_algorithms:performance:hardware}) and how to optimize a cryptographic algorithm implementation for a specific \gls{isa} (\Cref{sec:security_of_cryptographic_algorithms:performance:processors}). Finally, we discuss how profiling and benchmarking allows for identifying performance bottlenecks and consequently improve the overall efficiency of cryptographic algorithm implementation and execution (\Cref{sec:security_of_cryptographic_algorithms:performance:evaluation}).



\subsubsection{Efficiency Optimization Best Practices}\label{sec:security_of_cryptographic_algorithms:performance:best_practices}

We now discuss the most relevant high-level principles and methods for optimizing the efficiency of cryptography and cryptographic algorithm implementation:

\begin{itemize}
	
	\item \textbf{Cryptographic Algorithm Choice.} A basic aspect of efficiency optimization in cryptography is the choice of the most suitable cryptographic algorithm for the task at hand. In particular, when evaluating algorithms for data en/decryption, a fundamental distinction exists between symmetric and asymmetric algorithms:
	
	\begin{itemize}
		
		\item \textit{symmetric algorithms} such as AES and ChaCha use a single cryptographic key for both encryption and decryption of data. Symmetric algorithms are typically efficient and offer fast processing speeds thanks to their fairly simple mathematical structure. Usually, these algorithms involve computations like substitutions, permutations, XOR operations, and modular additions (over finite fields with order a power of 2) which processors can execute efficiently. Symmetric algorithms are ideal for applications where large volumes of data need to be en/decrypted such as in communications and secure storage;
		
		\item \textit{asymmetric algorithms} such as RSA and ECC use a pair of cryptographic keys, one public and one private, which act as a trapdoor to an underlying mathematical hard problem. Asymmetric algorithms are based on complex mathematical operations such as modular exponentiation, elliptic curve computations, and number-theoretic transforms. These operations introduce significant overhead, making asymmetric algorithms less suitable for bulk data encryption. Still, asymmetric algorithms play a crucial role in constructing various cryptographic algorithms, such as key encapsulation mechanisms and digital signatures.
		
	\end{itemize}

	We remark that, in practice, a hybrid approach is often employed to achieve the benefits of both types of encryption. 
	
	\item \textbf{Cryptographic Library Choice.} Once a cryptographic algorithm has been selected, chances are that an implementation of that algorithm within a cryptographic library already exists. In this case, the efficiency of such a library should be considered when deciding whether to re-implement the algorithm from scratch.
	
	\item \textbf{Memory Management.} Proper memory management can lead to significant efficiency improvements and is often a distinguishing factor between fast and slow implementations for cryptographic algorithms.
	
	\remarkBlock{Heap and stack}{The two primary types of memory allocation are stack and heap: the stack is usually used for local variables and function call data, while the heap is usually used for dynamic memory allocation. The stack is typically faster than the heap due to its simpler structure.}
	
	In particular, proper memory management expects to:
	
	\begin{itemize} 
		
		\item \textit{avoid (especially unintentional) data duplication} as it requires additional memory and incurs overhead, especially for large data structures --- such as those used in RSA or for elliptic curve points --- allocated in the heap. Instead, operations should be performed directly on the original data (\textit{in-place computations}) or using pointers to avoid unnecessary duplication;
		
		\item \textit{prefer arrays over vectors}, as arrays are typically stored as a contiguous block of memory in the stack, having a fixed size and fast access to data. Conversely, vectors tend to be stored in the heap, having a dynamic size and slow access to data. While offering greater flexibility compared to arrays, vectors come with a higher overhead, especially when their capacity is exceeded;
		
		\item \textit{employ suitable data structures} for organizing and accessing data efficiently. Examples of relevant data structures are:
		
		\begin{itemize} 
			
			\item \textit{linked lists}, which allow efficient insertion and deletion at both ends;
			
			\item \textit{hash maps}, which associate each element with a key for fast lookup;
			
			\item \textit{(binary) trees}, which allow efficient searching, insertion, and deletion operations.
			
		\end{itemize} 
		
		In general, complex arithmetic operations in cryptographic algorithm implementations may be more or less (memory) efficient according to how data are represented. For instance, a point on an elliptic curve is often stored as the $x$-coordinate only of an affine point $(x,y)$. However, computations performed on the projective plane $(x,y,z)$ are more efficient as they avoid the costly inversions typically needed in affine coordinates. On top of that, mixed operations between projective and affine coordinates allow for the direct use of large pre-computation tables, further enhancing performance;\footnote{A comprehensive list of formulae for each curve can be found in the Explicit-Formula Database (EFD) available at \url{https://www.hyperelliptic.org/EFD/}}
		
		\item \textit{rely on pre-computation}, that is, compute specific values (e.g., lookup tables such as S-boxes in AES) before the execution of cryptographic algorithm implementation. Intuitively, the availability of such values prevents redundant computations and reduces time complexity.
		
	\end{itemize}
	
\end{itemize}
	




\subsubsection{Languages and Compilers for Efficiency}\label{sec:security_of_cryptographic_algorithms:performance:languages_and_compilers}

Below, we investigate the impact that the choice of the programming language (\Cref{sec:security_of_cryptographic_algorithms:performance:languages_and_compilers:languages}) --- and, consequently, compiler (or interpreter) and configuration options (\Cref{sec:security_of_cryptographic_algorithms:performance:languages_and_compilers:compilers}) --- may potentially have on the efficiency of the implementation of cryptographic algorithms. In other words, we are concerned about whether the choice of the programming language and compiler or interpreter can improve or worsen efficiency.


\paragraph*{Programming Languages for Efficiency.}

Compiled languages (like C and C++, Go, and Rust) are generally more performant than both runtime-based (like Java, C\#, and Erlang) and interpreted languages (like Python, JavaScript, and Ruby). In fact, compiled languages produce machine code (see \Cref{sec:security_of_cryptographic_algorithms:performance:processors}), while runtime-based and interpreted languages have notable latency. Moreover, compiled languages allow for fine-tuned control over system resources. Below, we discuss the compiled languages among those already considered in \Cref{sec:security:languages_and_compilers:languages}:

\begin{itemize}
	
	\item \textbf{C and C++.} Despite ongoing efforts to replace them with newer languages, C and C++ have proven to be the preferred languages for performance-critical applications --- including cryptography. In fact, many widely used cryptographic libraries (e.g., OpenSSL) are written in C and C++, benefiting from years of continuous optimization and refinement. In detail, C and C++ offer minimal abstractions when compared to higher-level languages (e.g., Java), providing developers with fine-grained control over efficiency;
	
	\item \textbf{Rust.} Emerged as a modern alternative to C and C++ for performance-critical applications, Rust provides zero-cost abstractions, powerful concurrency features, and modern tooling (such as the Cargo package manager) enabling developers to write efficient and secure code. In fact, the adoption of Rust for cryptographic algorithm implementation is steadily growing. In terms of efficiency, the main drawback of Rust is slow compilation,\footnote{rust-lang (\url{https://prev.rust-lang.org/en-US/faq.html})} a direct result of the security features provided by its static analysis tools (see \Cref{sec:security:languages_and_compilers:languages});
	
	\item \textbf{Go.} Designed at Google for efficiency and ease of use. Go has many built-in tools for parallelizing computations. Go uses a garbage collector to manage memory which, although simplifying software development, hinders performance. Still, Go is often seen as an acceptable choice to develop scalable and efficient applications quickly, especially for back-end software and security-related systems.
	
\end{itemize}

A comprehensive and detailed analysis of the efficiency of each of the aforementioned languages (and many more) can be found in \cite{pereira_ranking_2021}.

Besides the aforementioned languages, another language (or, better, type of languages) relevant to efficiency is \textbf{assembly}, which allows for even more precise control over efficiency --- for more details, see Note 3.4.2 in \Cref{sec:security_of_cryptographic_algorithms:performance:processors}. Finally, we also mention \textbf{WebAssembly} (also called Wasm), which is a standard describing a low-level binary instruction format to run high-performance applications across any platform with a web browser. In brief, WebAssembly allows developers to write parts of their program in languages such as C, C++, or Rust. Such parts are then compiled into WebAssembly bytecode and executed within web browsers with little extra overhead. Although WebAssembly bytecode is not as fast as running native code directly on the machine, it provides a way to get near-native performance in the browser, making it a great option for web-based cryptographic applications. Additionally, WebAssembly works in a secure, isolated environment, which helps protect cryptographic code. In other words, WebAssembly is particularly relevant to web applications (at client-side).


\paragraph*{Compiler and Interpreters Options for Efficiency.}

Compiler optimization flags are used to improve the efficiency of software by altering how the compiler generates machine code. A general rule of thumb is that, while compiler optimizations can improve the runtime efficiency of a program, they may increase the compile time and the binary size. These optimizations are typically achieved through techniques such as constant folding, dead code elimination, and aggressive optimizations including:

\begin{itemize}
	
	\item \textbf{inlining}, which replaces function calls with the actual code of the function, eliminating the overhead of function calls;
	
	\item \textbf{vectorization}, which converts loops and operations to use \gls{simd} instructions (see \Cref{sec:security_of_cryptographic_algorithms:performance:processors} for more details on \gls{simd} instructions);
	
	\item \textbf{\gls{lto}}, which optimizes the program across multiple translation units during the linking stage, allowing the compiler to perform optimizations that require access to the whole code of the program. \gls{lto} can lead to significant performance improvements, as it allows the compiler to apply optimizations that are not possible when each source file is compiled independently.
	
\end{itemize}

The aforementioned are generic optimizations that almost all compilers aim to achieve. However, specific compilers may have specific flags enabling further optimization techniques. For instance:

\begin{itemize}
	
	\item \textbf{GCC and Clang}, which are popular compilers for C and C++, allow basic optimization with the flag \texttt{-O}, progressively increasing through \texttt{-O2} and \texttt{-O3}, with the latter being the most aggressive. In GCC and Clang, \gls{lto} can be enabled with the \texttt{-flto} flag. A comprehensive list of all optimizations flags can be found in the GNU Compiler Collection Flags\footnote{GNU Compiler Collection Flags (\url{https://www.spec.org/cpu2017/flags/gcc.2018-02-16.html})} and in the Clang documentation;\footnote{Clang 20.0.0git documentation (\url{https://clang.llvm.org/docs/CommandGuide/clang.html})}
	
	\item \textbf{rustc}, which is the compiler of Rust, allows for using optimization flags even though it is usually more common to rely on Cargo (the build system of Rust) for managing optimizations. Rust uses the \texttt{-{}-release} flag to apply the optimizations specified in the \texttt{[profile.release]} section of the \texttt{Cargo.toml} file. Similar to GCC and Clang, there is an \texttt{opt-level} setting that ranges from 1 to 3, with higher levels applying progressively more aggressive optimization strategies. Additionally, setting \texttt{lto = true} enables \gls{lto}. For direct use of \texttt{rustc}, it is possible to specify optimizations with the \texttt{-{}-opt-level} flag, which works similarly to \texttt{cargo}'s configuration. A list of possible optimizations can be found in the rustc documentation\footnote{The rustc book (\url{https://doc.rust-lang.org/rustc/codegen-options/index.html?highlight=opt-level})}.
	
\end{itemize}



\subsubsection{Hardware for Efficiency}\label{sec:security_of_cryptographic_algorithms:performance:hardware}

Hardware-based cryptographic modules (see CSC-D1---§Section 2) and cryptographic accelerators include specialized circuits and instruction sets that significantly (i.e., usually at least one order of magnitude \cite{baldanzi_crypto_2019}) improve performance \cite{meurice_de_dormale_high_2007} and energy efficiency \cite{haas_evaluating_2013} during the execution of cryptographic algorithms --- e.g., by accelerating or parallelizing particular operations. Typically, the more a piece of hardware is tailored for specific operations, the better its performance and energy efficiency. In fact, cryptographic accelerators are often designed for specific cryptographic algorithms (e.g., AES) and are deployed as hardware side components, integrated circuits, or co-processors near general-purpose hardware.\footnote{General-purpose hardware includes both processors (i.e., CPUs) and GPUs. Note that general-purpose GPUs are sometimes used for accelerating cryptographic computations, having often better performance than general-purpose processors in executing arithmetic operations} The main types of cryptographic accelerators include:
\begin{itemize}
	\item \textbf{\glspl{asic}}, which are designed for executing specific (sets of) operations, usually supporting hash functions (e.g., in cryptocurrency mining) and en/decryption --- but also other operations such as random numbers generation. \glspl{asic} are not reconfigurable and cannot be supplied with additional instructions after design and manufacturing. The two primary design methods for \glspl{asic} are gate-array design and full-custom design: the former uses predefined components to minimize the design effort and associated costs, while the latter provides better flexibility and performance at the cost of a more complex and costly design. Examples of \glspl{asic} are Intel QuickAssist Technology (QAT) \glspl{asic}\footnote{Intel® QuickAssist Technology (Intel® QAT) (\url{https://www.intel.com/content/www/us/en/architecture-and-technology/intel-quick-assist-technology-overview.html})} and ARM CryptoCell;\footnote{CryptoCell-300: Embedded Security for High-Efficiency SoCs – Arm®  (\url{https://www.arm.com/products/silicon-ip-security/crypto-cell-300})} 
	
	\item \textbf{\glspl{fpga}}, which comprise programmable logic blocks (i.e., logic gates, flip-flops, look-up tables) and programmable interconnects (i.e., pathways for signals to travel between logic blocks) allowing for reconfigurable hardware operations. In other words, \glspl{fpga} are not application-specific and can be repeatedly re-programmed through hardware description languages (such as Verilog\footnote{IEEE SA - IEEE 1364-2001 (\url{https://standards.ieee.org/ieee/1364/2052/})} and VHDL\footnote{IEEE SA - IEEE 1076-2019 (\url{https://standards.ieee.org/ieee/1076/5179/})}) to suit different applications after design manufacturing --- as opposed to \glspl{asic}. Hence, \glspl{fpga} bridge the gap between general-purpose hardware and fully specialized hardware like \glspl{asic}. Consequently, \glspl{fpga} offer more performance but less flexibility than general-purpose hardware, and more flexibility but less performance than \gls{asic}. Examples of \glspl{fpga} are Xilinx \glspl{fpga} by AMD\footnote{AMD FPGAs (\url{https://www.amd.com/en/products/adaptive-socs-and-fpgas/fpga.html})} and Intel Stratix 10 series;\footnote{Intel® Stratix® 10 FPGAs Overview - High Performance Intel® FPGA (\url{https://www.intel.com/content/www/us/en/products/details/fpga/stratix/10.html})}
	
	\item \textbf{specialized instruction sets}, which are integrated into the design of general-purpose hardware but require dedicated hardware within the processors themselves (i.e., not on a separate co-processor) \cite{cheng_sok_2024}. Examples of specialized instruction sets are Intel AES-NI\footnote{Intel® Advanced Encryption Standard Instructions (AES-NI) (\url{https://www.intel.com/content/www/us/en/developer/articles/technical/advanced-encryption-standard-instructions-aes-ni.html})} and ARMv8 Cryptographic Extension.\footnote{Arm® Cortex-A78 Core Cryptographic Extension Technical Reference Manual (\url{https://developer.arm.com/documentation/101432/r1p2/Functional-description/About-the-Cryptographic-Extension})}
	
\end{itemize}

We highlight that, besides efficiency, hardware-based cryptographic modules provide further benefits such as security (see \Cref{sec:security:hardware}) and key management --- which many but not all cryptographic accelerators provide. As a final remark, we note that the need for hardware support increases when considering resource-intensive advanced cryptosystems (e.g., homomorphic encryption) and post-quantum cryptographic algorithms.



\subsubsection{Optimizing for Instruction Set Architectures}\label{sec:security_of_cryptographic_algorithms:performance:processors}

An \gls{isa} is often defined as an abstract model\footnote{What is Instruction Set Architecture (ISA)? – Arm® (\url{https://www.arm.com/glossary/isa})}$^,$\footnote{Instruction Set Architecture (\url{https://www.intel.com/content/www/us/en/developer/tools/isa-extensions/overview.html})} or interface between software and hardware describing the capabilities of one or more processors and how to control them. In other words, an \gls{isa} can be viewed as a developers' manual and standard to interact with the processors. More in detail, an \gls{isa} defines supported data types, registers, memory management, key features (e.g., virtual memory), input/output model and, especially, instructions and their binary encoding. Being an abstract model, there usually exist multiple implementations (in this case, processors) of an \gls{isa} with different characteristics --- such as performance, physical size, and monetary cost --- capable of running the same instructions. Hence, \glspl{isa} ensure compatibility and portability across different processors. An \gls{isa} can also be extended with further instructions to exploit new hardware (e.g., see \Cref{sec:security_of_cryptographic_algorithms:performance:hardware}), improve performance, and satisfy the needs of a specific application or domain; such extensions are typically backward-compatible. Intuitively, knowing and understanding an \gls{isa} allows developers to write and compile more efficient code. \glspl{isa} are commonly categorized according to how instructions correspond to arithmetic and memory access operations:

\begin{itemize}
	
	\item \textbf{\gls{risc}}, which comprises \glspl{isa} where instructions correspond to single low-level operations (e.g., load data from memory, perform an arithmetic operation, store data in memory). As a consequence, \gls{risc} \glspl{isa} typically prioritize simplicity, power efficiency, and performance by expecting (optimized and specialized) instructions with uniform length and predictable execution time, making instruction pipelining easier. Generally, \gls{risc} \glspl{isa} include frequently used instructions (i.e., operations) only, while less common operations need to be implemented as subroutines, resulting in a small (but inherently infrequent) overhead.
	
	\remarkBlock{The meaning of ``reduced'' in \gls{risc} \glspl{isa}}{The term ``reduced'' in ``\glsfirst{risc}'' implies that each instruction corresponds to a small and predictable amount of work (e.g., a single data memory cycle) and not that the set of instructions itself is small --- although it is usually the case.}
	
	Examples of \gls{risc} \glspl{isa} are those of ARM (e.g., Armv8-A for 64-bit processors only, Armv7-A for 32-bit processors only, Thumb for 16-bit processors only) --- common in mobile and embedded devices --- and RISC-V;
	
	\item \textbf{\gls{cisc}}, which comprises \glspl{isa} where single instructions are capable of multi-step operations and can correspond to several low-level operations (e.g., load data from memory, perform an arithmetic operation, and also store the result). Consequently, \gls{cisc} \glspl{isa} typically prioritize flexibility, usability, abstraction, fewer main memory accesses, and compact code. Examples of \gls{cisc} \glspl{isa} are those of Intel and AMD (e.g., x86-64/AMD64 for processors up to 64-bit, x86 for processors up to 32-bit) --- common in servers and desktops.
	
\end{itemize}

Finally, we note that some consider further categories such as \gls{misc} and \gls{oisc} and that there exist special \glspl{isa} such as \gls{epic} and \gls{vliw} used for particular applications or domains.

\remarkBlock{\gls{isa} vs. assembly}{Processors execute a stream of binary-encoded instructions called ``machine code'' for which there usually exist a more human-readable rendition called ``assembly code''. The term ``assembly'' identifies a low-level programming language that uses operation codes (or ``opcodes'') --- represented with mnemonics --- to enumerate instructions and symbolic names to refer to registers. Assembly code has an almost direct mapping to machine code, although an assembly language may comprise some higher-level syntactical constructions and allow for inserting developers' comments. Intuitively, there does not exist a single assembly language, but rather multiple assembly languages, each of which corresponds to a specific \gls{isa} --- while each \gls{isa} may have one or more assembly languages; for instance, x86 processors have two assembly languages, the AT\&T syntax (also called GAS) and the Intel syntax (also called NASM). Finally, the translation from assembly code to machine code is done by the assembler.}

One of the simplest and most effective ways to optimize cryptographic algorithms for specific \glspl{isa} is to leverage compilers and interpreters options (see \Cref{sec:security_of_cryptographic_algorithms:performance:languages_and_compilers}) to specify target processors and \glspl{isa} and optimization levels. For instance, the \texttt{gcc} compiler for \texttt{C} has a flag \texttt{march} which, when set to \texttt{native} (i.e., \texttt{-march=native}), asks the compiler to fine-tune the compiled code for the \gls{isa} of the host processor. Similarly, the \texttt{-ffast-math} flag allows the compiler to diverge from the IEEE 754 standard \cite{ieee_754} for floating point operations and improve their performance according to the underlying \gls{isa} (at the cost of accuracy). In general, it is important to rely on compilers and interpreters which consider code optimization for the underlying \gls{isa}; intuitively, this requires deep knowledge of such compilers and interpreters.

Another way to optimize the execution of cryptographic algorithms is to study specific software optimization guides,\footnote{Software optimization resources. C++ and assembly. Windows, Linux, BSD, Mac OS X (\url{https://www.agner.org/optimize/})} official documentation and manuals,\footnote{Intel® 64 and IA-32 Architectures Optimization Reference Manual Volume 1 (\url{https://www.intel.com/content/www/us/en/content-details/814198/intel-64-and-ia-32-architectures-optimization-reference-manual-volume-1.html})} and leverage specialized instructions and extensions (e.g., Intel AES-NI; see \Cref{sec:security_of_cryptographic_algorithms:performance:hardware}) of the underlying \gls{isa}; intuitively, this requires deep knowledge of the \gls{isa}. An instance of such specialized instructions are the \gls{simd} instructions which --- by suitably placing distinct small pieces of data (usually, vectors) into the same registers --- perform the same (e.g., arithmetic) operation on each piece of data in parallel, yielding an almost proportional increase in performance (net of downclocking); this practice is often called vectorization. For instance, ARM \glspl{isa} have the \texttt{NEON} extension which groups \gls{simd} instructions operating on vectors; Intel \glspl{isa} have a similar extension called \texttt{AVX2}. Specialized instructions can be leveraged either automatically by compilers and interpreters --- e.g., with the \texttt{pragma} directive in \texttt{gcc} or by specialized compilers such as the \gls{ispc}\footnote{Intel® Implicit SPMD Program Compiler (\url{https://ispc.github.io/})} --- or manually by developers which can use:

\begin{itemize}
	
	\item \textbf{compiler intrinsics}, which allows for accessing \gls{isa} instructions without delving into assembly code, and are thus a convenient way to exploit specialized instructions and \gls{isa} extensions in high-level code. In brief, intrinsic allows for using functions which are high-level language representations of machine instructions. For example, in C++, it is possible to use an intrinsic function that corresponds to a vector addition machine instruction, eliminating the need for inlining assembly code. With intrinsics, the compiler handles most of the optimizations, enabling the developers to focus on algorithm design and data organization;
	
	\item \textbf{inline assembly}, which is embedded directly in the program, allowing the developers to fine-tune specific parts of the software near the high-level language. Inline assembly is useful in presence of particular instructions that the compiler may not handle efficiently on its own, or when optimizations based on a specific \gls{isa} are needed;
	
	\item \textbf{external assembly}, which is assembly code written in separate files with respect to the rest of the program. External assembly is the most flexible but also requires the most effort and expertise. However, note that, when using external assembly, the compiler may not be able to fully optimize the surrounding high-level code, as the external assembly code is often opaque to the compiler. As a result, programmers may lose the benefits of some automatic optimizations, which can ultimately negate, if not worsen, the performance gain.
	
\end{itemize}

Intuitively, all of the aforementioned three approaches require deep knowledge of the assembly language. It is worth noticing that there may already exist libraries implementing arithmetic operations and mathematical routines (e.g., linear algebra\footnote{OpenBLAS (\url{http://www.openmathlib.org/OpenBLAS/})}) strongly optimized for a specific \gls{isa}.

Finally, the efficiency of cryptographic algorithms can be improved by optimizing memory access and alignment, pre-fetching (to reduce cache misses), and minimizing cache contention (to reduce cache thrashing) to decrease memory latency and maximize data throughput. In this regard, it is particularly important to avoid data hazards, i.e., situations where the execution of an instruction relies on the outcome of a preceding instruction that is still in progress. Memory access optimization requires deep knowledge of the latency and throughput of each instruction,\footnote{Latency and throughput for a specific \gls{isa} are usually described in instruction tables such as those available in \url{https://www.agner.org/optimize/instruction_tables.pdf}} their interaction with other instructions, the pipeline functioning, and the memory management of the \gls{isa}. 

\remarkBlock{The importance of profiling and benchmarking}{Profiling and Benchmarking --- which are discussed in \Cref{sec:security_of_cryptographic_algorithms:performance:evaluation} --- are of the utmost importance to identify performance bottlenecks and determine whether certain compiler and interpreter options, \gls{simd} instructions, or memory access optimization techniques actually improve the efficiency of cryptographic algorithm implementation.}

To conclude, we remark that --- besides the \gls{isa} abstract model --- the efficiency of cryptographic algorithm implementation heavily depends on the characteristics of the underlying processor (clock frequency, multi-threading, support for parallel and distributed computing, instruction pipelining, out-of-order execution, speculative execution, dedicated hardware for supporting \gls{isa} extensions). 



\subsubsection{Evaluation and Benchmarking of Efficiency}\label{sec:security_of_cryptographic_algorithms:performance:evaluation}

Benchmarking the efficiency of cryptographic algorithm implementation offers several benefits. In particular, benchmarking allows for comparing different implementation approaches and, more importantly, helps identify performance bottlenecks, revealing what portions of the code need to be refined.


\paragraph*{Efficiency Metrics for Comparing cryptographic algorithm implementation.}

When benchmarking cryptographic algorithm implementation, there are several important metrics to consider. Intuitively, the relevance of those metrics may vary according to the underlying application and domain in which cryptographic algorithms are employed. The most important metrics that are often collected through benchmarking are:\footnote{eBACS: ECRYPT Benchmarking of Cryptographic Systems (\url{https://bench.cr.yp.to/})})

\begin{itemize}
	
	\item \textbf{latency}, which is the (time) interval between the start and completion of a cryptographic computation. Latency is usually measured either in time (e.g., milliseconds) or processor cycles per operation. Both measurements are valuable: time provides real-world user-friendly benchmark results and comprises factors such as processor speed, memory availability, and other computational resources; processor cycles focus on the number of instructions executed by the processor, which is useful for low-level optimization. Latency is typically the most important metric for asymmetric cryptography, as the amount of data to process is usually small and fixed;
	
	\item \textbf{throughput}, which is the rate at which data are elaborated. Throughput is typically measured as operations per second (ops/sec) or processor cycles per byte (cycles/byte). Throughput is particularly useful when the size of the input data varies significantly. Also, throughput is especially relevant to symmetric cryptography and hashing, where large amounts of data may need to be elaborated;
	
	\item \textbf{initialization}, which refers to the setup phase of a cryptographic algorithm before executing its first cryptographic computation. Initialization is measured similarly to latency, i.e., in time or processor cycles. Concretely, initialization often involves tasks such as key generation, loading and parsing of configuration parameters, and data structure initialization. Although initialization is usually a one-time cost, its relevance increases when re-keying is frequently required (e.g., during a TLS handshake) and it is a key factor in determining key agility;
	
	\item \textbf{size}, which refers to the length --- or, more generally, dimensions --- of private, public, and secret keys, as well as the size of the output of cryptographic computations (e.g., ciphertexts, digital signatures). Size is an important metric for evaluating the efficiency of cryptographic algorithm implementation in terms of memory usage and network bandwidth. Hence, size is particularly relevant when information needs to be transmitted or stored for a long time;
	
	\item \textbf{scalability}. which measures the efficiency of cryptographic algorithm implementation under varying parameters. Scalability includes testing efficiency with different data sizes (e.g., larger payloads to encrypt) and varying levels of security (e.g., key size), which typically results in some sort of performance degradation. Scalability is an essential metric for assessing how well cryptographic algorithm implementation performs as the workload increases.
	
\end{itemize}

Measuring the aforementioned metrics (and making results public\footnote{SUPERCOP (\url{https://bench.cr.yp.to/supercop.html})}) allows for determining which cryptographic algorithm and implementation is the most suitable according to the efficiency requirements of the considered application and domain. Additionally, benchmarking helps in identifying unexpected slowdowns and performance bottlenecks, guiding advanced analysis and fine-tuned improvements. Intuitively, benchmarking is the most useful when measurements are accurate, which is not trivial and instead requires considering several challenges. In fact, performance on a specific machine (e.g., server, smartphone) can vary significantly due to, e.g., process scheduling, concurrency, and caching. In other words, simply measuring the time required for a specific function to complete is not often enough to produce reliable benchmarking results. Additionally, cryptographic computations often involve complex arithmetic, and the time taken may vary considerably depending on the specific key used, even when the key size is the same.

To address these challenges and promote accurate measurements, several strategies can be taken. Using a large sample size, implementing a warm-up phase, and repeating benchmarks multiple times are all strategies commonly adopted for reducing statistical noise in the measurements and averaging the results obtained. However, we note that an even more effective strategy consists of relying on benchmarking frameworks that already handle these challenges and implement the aforementioned strategies. Examples of well-established benchmarking frameworks are the Google benchmark\footnote{Google benchmark (\url{https://github.com/google/benchmark})} for C and C++, Criterion-rs\footnote{Criterion-rs (\url{https://bheisler.github.io/criterion.rs/book/})} for Rust, and the standard library\footnote{Go standard-library (\url{https://pkg.go.dev/testing})} for Go.


\paragraph*{Efficiency Metrics for Resource Utilization.}

When optimizing the implementation of cryptographic algorithms, it is often useful to focus on resource utilization, i.e., how efficiently the implementation uses available resources such as processor, memory, and cache:

\begin{itemize}
	
	\item \textbf{processor usage}, that is, in which portions of the code the processor is spending the most time. Analyzing processor usage can identify which code paths consume excessive processing capacity and steer optimization efforts accordingly;
	
	\item \textbf{cache misses}, that is, data required by the process is not available in the processor cache and need therefore to be fetched from slower memory (e.g., the main memory). Cache misses can introduce significant performance bottlenecks, as accessing the main memory is much slower than accessing the processor cache. Hence, monitoring for cache misses can help in identifying inefficient memory access patterns;
	
	\item \textbf{memory footprint}, which is the overall amount of memory used by the process, including keys, buffers, and temporary variables. Memory footprint is relevant as, in some contexts, memory consumption might be even more critical than raw performance, particularly in resource-constrained devices (e.g., embedded systems), where limited memory availability can directly affect performance and stability.
	
\end{itemize}

\paragraph*{Profiling.} Resource utilization is often analyzed through profiling, a form of dynamic analysis aimed at gathering process events. Profilers can be broadly divided into two categories:

\begin{itemize} 
	
	\item \textbf{statistical profilers} which operate by sampling the process call stack at regular intervals using system interrupts. The data collected provides a statistical approximation and can offer an accurate depiction of the performance of the process. An example of a statistical profiler is the Linux perf tool;\footnote{perf (\url{https://www.brendangregg.com/overview.html})}
	
	\item \textbf{instrumentation-based profilers} which inject instructions into the target process to gather process events. These profilers allow for obtaining precise data but may themselves alter the performance of the process due to the overhead introduced by the injected instructions. An example of an instrumentation-based profiler is Valgrind;\footnote{Valgrind (\url{https://valgrind.org/})}
	
\end{itemize}

Usually, statistical profilers are used first while instrumentation-based profilers are used for very precise and targeted analysis only. In both cases, the data gathered by a profiler is often difficult to interpret. Some visualization tools that can help are flame-graphs\footnote{flame-graphs (\url{https://www.brendangregg.com/overview.html})} and KCachegrind.\footnote{KCachegrind (\url{https://kcachegrind.sourceforge.net/html/Home.html})}
